\subsection{Purpose}
Internships are a vital bridge between academic learning and professional experience, providing students with opportunities to apply their skills and companies with access to emerging talent. However, finding the right match between students and internships can be a complex and inefficient process.
The Students\&Companies (S\&C) platform addresses these challenges by providing an integrated system that connects university students with companies offering internships. S\&C provides a wide range of features offering a seamless experience for students from the search of an internship to its completion and for companies an easy way to advertise their internship projects. The platform also lets universities to monitor and support their students throughout their internship, overseeing its execution and supporting the involved parties, should problems arise.
\subsubsection{Goals}
\begin{goals}
\item Students must be able to search for internships posted by companies.
\item Companies must be able to upload their internships offers.
\item Students must be notified when an internship that could interest them becomes available.
\item Companies must be notified about available students with CVs and skills matching their internship requirements.
\item Students must be able to apply for internships of interest.
\item Companies must be able to accept applicants to their internship.
\item Companies must be supported in conducting a selection process for applicants, including the collection of information about the students using questionnaires and the management of interviews.
\item Students want to receive suggestions on how to make their CVs more appealing.
\item Companies want to receive suggestions on how to make their project descriptions more appealing.
\addtocounter{goalsi}{1}
\goalitem{A} Students must be able to keep track and monitor the execution and the outcomes of the matchmaking process.
\goalitem{B} Companies must be able to keep track and monitor the execution and the outcomes of the matchmaking process.
\goalitem{C} Universities must be able to monitor the status of the internship in which its students have applied.
\item The interested parties of an internship (students, companies and universities) must be able to communicate with each other about info and eventual problems regarding the ongoing internship.
\addtocounter{goalsi}{1}
\goalitem{A} Companies must be able to file complaints about students interning with them.
\goalitem{B} Students must be able to file complaints about the companies where they are interning.
\item Universities must be able to review complaints about their students and the companies where their students are interning.
\item Universities must be able to interrupt internships between their students and a company.
\end{goals}
\subsection{Scope}
The Students\&Companies platform provides its services to different users, identifiable into three main classes: students,  companies and universities. Each one of them has different roles and interacts with the system in various ways.
S\&C is aimed mainly at university students that want to take part in an internship program, as they can search for new internships suiting their preferences or let the system inform them when internships matching their profile are available. S\&C is also able to provide suggestions on how to improve their profile to make it more appealing to companies.
The platform also supports companies easing the internships’ management. In doing so S\&C can be used to create questionnaires to select the most suitable candidates for an internship. Similarly to the suggestions for students, the system can analyze the internship offer to identify eventual improvements that can be made.
Hopefully, every internship comes to an end without problems but if that is not the case, interns and companies can communicate through the platform possibly involving the intern’s alma mater as a supervisor which can also terminate the internship if the parties cannot resolve the controversy.
\subsubsection{World Phenomena}
\begin{w-phenomena}
    \item Students want to do an internship.
    \item Companies offer an internship.
\end{w-phenomena}
\subsubsection{Shared Phenomena}
\paragraph{World Controlled}
\begin{s-phenomena}
    \item Companies upload their internship offers.
    \item Students search for an internship.
    \item Students accept the internship.
    \item Companies approve the internship contract for a student.
    \item Companies interview the student.
    \item Students give feedback on the internship.
    \item Companies give feedback on the internship.
    \item Students file complaints about the internship.
    \item Companies file complaints about the internship.
    \item University monitors the situation of internships.
    \item University receive a complaint.
    \item University responds to a complaint.
    \item University terminates an internship.
    \item Students keeps track, monitors the execution and the outcomes of the matchmaking process.
    \item Companies keeps track, monitors the execution and the outcomes of the matchmaking process.
\end{s-phenomena}
\paragraph{Machine Controlled}
\begin{s-phenomena}[resume]
    \item System informs the student when an internship that might interest him becomes available.
    \item System informs the companies about the availability of student CV corresponding to their needs.
    \item System establishes a contact between the student and the companies.
    \item System analyzes internship data for recommendations.
    \item System sends suggestions on how to improve the students CVs.
    \item System sends suggestions on how to improve companies project descriptions.
    \item System sends universities complaints.
    \item System sends a notification to the university about the complaint.
\end{s-phenomena}
\subsection{Definitions, Acronyms, Abbreviations}

\subsubsection{Definitions}
\begin{itemize}
    \item \textbf{Recommendations}: Process in which the system informs students when an internship that might interest them becomes available and can inform companies about the availability of student CVs corresponding to their needs.
    \item \textbf{SheerID}\cite{sheerid}: Platform that offers an API to verify students by using their academic credentials.
    \item \textbf{Contract}: The contract represents the lifecycle of an internship application, from initiation to its conclusion.
\end{itemize}
\subsubsection{Acronyms}
\begin{itemize}
    \item \textbf{S\&C}: Students \& Companies
\end{itemize}
\subsubsection{Abbreviations}
\begin{itemize}
    \item \textbf{G*}: goals.
    \item \textbf{WP*}: world phenomena.
    \item \textbf{SP*}: shared phenomena.
    \item \textbf{R*}: functional requirements.
    \item \textbf{UC*}: use cases.
\end{itemize}
\subsection{Revision History}
\begin{itemize}
    \item 22/12/2024 - Version 1.0.
\end{itemize}
\subsection{Reference Documents}
\begin{itemize}
    \item Assignment document “Assignment RDD AY 2024-2025.
\end{itemize}
\subsection{Document Structure}
\textbf{Introduction}: This section provides a concise overview of the project, focusing on the objectives and the reasons for its development. It outlines the key goals that the project aims to achieve.
\textbf{Overall description}: This section offers a high-level explanation of how the system operates, providing a detailed account of the phenomena involved in the interaction between the world, the machine, or both. It also includes a description of the domain, along with the assumptions that underpin the system’s design and functionality. \newline
\textbf{Specific requirements}: Here, a thorough analysis of the requirements necessary to meet the project’s goals is provided. This section also includes important information for developers, such as hardware and software constraints, along with any other specifications that guide the development process. \newline
\textbf{Formal analysis using Alloy}: This section presents a formal representation of world phenomena, using the Alloy language to describe the system’s behavior and properties in a precise and logical manner. \newline
\textbf{Effort spent}: This section details the time and resources invested in the creation of the document, broken down by each section for clarity and transparency.
References: This section includes citations and references to any documents, frameworks, or software tools utilized in the development of this paper.